\section{Conclusion}

Tool-agent safety requires a monitor to mediate a conservative abstraction of the effects a call will commit, not merely its name or surface form. Counterfactual onboarding intervenes on calls and state to test whether candidate atoms preserve observed security-effect distinctions; deterministic pre-commit checks then bind the frozen abstraction to task-scoped authority. The resulting confinement argument is useful precisely because it exposes causal-abstraction, authority, mediation, check--use, and uncertainty obligations that aggregate attack success can hide. Our current prototype establishes exact sandbox interception and exposes concrete contract, envelope, interaction, and default-semantics gaps; it does not yet discharge the complete architecture. Common-protocol baselines, long-trajectory tests, causal ablations, and interface-cost measurements determine whether the hardened construction improves security without merely suppressing useful execution.
